
% Color del resalte trasparentado a la derecha del índice
\begin{tikzpicture}[remember picture, overlay]
    \draw[line width=0.73\textwidth,color=Resalte!65!white] ($(current page.north east) - (0,0)$) -- ($(current page.south east) - (0,0)$);
\end{tikzpicture}
% Columna de la izqueirda, está el índice y el Sobre Momentum:
\begin{minipage}[t]{0.72\textwidth}
    \vspace{-0.3cm}
    \begin{minipage}[t]{\hsize}
        \makebox[0pt][l]{
            % Aquí incluimos el tableofcontens,
            \begin{adjustbox}{valign=t,minipage={0.9\textwidth},margin={0pt,0pt}}
                {   % Indice personalizable, con tamaño, color, fuente...
                    \section*{\Huge \textbf{Índice}}
                    \begingroup
                    % Esto es para que no salte el Indice default, y lo podamos poner nosotros
                    \let\contentsname\relax
                    \tableofcontents
                    \endgroup
                }%
            \end{adjustbox}%
        }%
        % Texto debajo del índice, sobre momentum
        \vfill
        \begin{minipage}[t]{0.9\textwidth}
            \vfill
            \vspace{2cm}
            {\Huge \color{Resalte!75!black} \textbf{Sobre Momentum}:} \\[1cm]Tras uns cuantos meses de traballo, por fín podemos dar saída á nova revista estudiantil \textit{momentum}, unha revista que busca ser un medio de comunicación tanto dentro coma fóra da facultade de física onde expor os intereses científicos do estudantado. Cuestións de Divulgación, Actualidade, Entrevistas, Historia, Filosofía da Ciencia, Opinión, Programación... Son todos temas que teñen cabida dentro deste proxecto. Esta revista está realizada integramente polo estudantado da Facultade de Física USC, onde pretendemos ter un recuncho de expresión máis aló do estrictamente académico. Fundada no ano 2025 co obxectivo de persistir na historia, esforzámonos sempre en mellorar. Non dubidedes en deixar a vosa pegada!
        \end{minipage}
    \end{minipage}
\end{minipage}
\hfill
% Texto que está en la linea roja trasparente, logos y autores/editores/directores
\begin{minipage}[t]{0.28\textwidth}
    \vspace{-0.3cm}
    \hspace{-0.7cm}
    \begin{center}
        \textcolor{white}{\Large
        \textbf{\today }\\[5mm]
        \textbf{Número 001} \\[1.5cm]
        }
    \end{center} 
    \Autores \\[1.5cm]
    \begin{tikzpicture}[scale=1.5]
       % \draw[thick, white] (0,0.5) -- (2.85,0.5);
        \node at (0,0) {\includegraphics[width=0.7cm]{logos/GitHub-branco.png}};   
        \begingroup
        \hypersetup{urlcolor=white}
        \node[anchor=west,align=left] at (0.5,0) {\textbf{\textcolor{white}{\href{https://github.com/DAF-USC/revista}{GitHub}
        }}};
        \endgroup
        \node at (0,-1) {\includegraphics[width=0.7cm]{logos/Icone_Email.png}};
        \node[anchor=west,align=left] at (0.5,-1) {\footnotesize \textbf{\textcolor{white}{revistafisicausc@gmail.com}
        }};
     %   \draw[thick, white] (0,-1.5) -- (2.58,-1.5);
    \end{tikzpicture} \\[1.7cm]
    \begin{center}
    \includegraphics[width=0.8\linewidth]{./logos/usc-invertido.png} % Logo 2
    \end{center}
\end{minipage}
